% Le petit * enlève le numéro (1.)
\section*{Introduction Générale} 
% Cette ligne force LaTeX à ajouter l'intro dans le sommaire même si elle n'a pas de numéro
\addcontentsline{toc}{section}{Introduction Générale}

Le secteur du e-commerce est marqué par une concurrence accrue où la fidélisation des clients est devenue un enjeu stratégique majeur. Dans ce contexte, ce projet porte sur l'analyse et la prédiction du comportement des clients d'une boutique de vente de cadeaux en ligne.

L'objectif principal est de développer un modèle de Machine Learning capable d'anticiper le \textbf{Churn} (l'attrition), c'est-à-dire d'identifier les clients qui risquent de ne plus effectuer d'achats sur la plateforme.

\begin{tcolorbox}[title=Objectifs du Projet]
	\begin{itemize}
		\item \textbf{Exploration} : Comprendre les habitudes d'achat et les profils types des clients.
		\item \textbf{Prédiction} : Anticiper le départ des clients grâce à des algorithmes de classification.
		\item \textbf{Actionnabilité} : Fournir des indicateurs permettant à l'entreprise de mettre en place des campagnes de rétention ciblées.
	\end{itemize}
\end{tcolorbox}

% L'étoile enlève aussi le numéro de la sous-section
\subsection*{Présentation du Dataset}
\addcontentsline{toc}{subsection}{Présentation du Dataset}

La base de données utilisée (\textit{retail\_customers\_COMPLETE\_CATEGORICAL.xlsx}) regroupe des informations riches et variées sur le parcours client, notamment :
\begin{itemize}
	\item \textbf{Variables RFM} : Récence, Fréquence et Montant total des achats.
	\item \textbf{Données Démographiques} : Âge, pays de résidence, genre.
	\item \textbf{Comportement Digital} : Temps de connexion préféré, saison favorite, nombre de tickets de support ouverts.
	\item \textbf{Variable Cible} : La colonne \texttt{Churn} (0 pour fidèle, 1 pour parti).
\end{itemize}

Ce rapport détaille la démarche scientifique suivie, de la préparation rigoureuse des données jusqu'au déploiement d'une solution prédictive.